% 调用上级目录公共模板
\documentclass{../template/softeng}

\begin{document}
	
% 封面：文档名、项目名、版本
\doccover{软件需求规格说明书}
{在线论文自动生成系统}
{V2026.07}
	
% 目录
\tableofcontents
\newpage
	% 修订记录
\begin{revrecord}
	V1.0 & 2026-07-13 & 郑羽婷 & 陆思雨 & 完成功能性需求 \\
	\hline
	V1.1 & 2026-07-14 & 陆思雨 & 郑羽婷 & 扩充需求描述 \\
	\hline
	V1.2 & 2026-07-14 & 钟晶卉 & 陆思雨 & 扩充需求描述\\
	\hline
	V1.2 & 2026-07-16 & 陆思雨 & XXX & 修改掉了技术实现相关，扩充更多业务相关，完成最终版\\
	\hline
\end{revrecord}
	
% 正文结构（软件工程国标标准章节）
\section{引言}
\subsection{编写目的}
本文档为"在线论文文档自动生成工具"软件需求规格说明书，用于完整定义系统面向用户的功能需求与业务规则，明确项目边界与使用约束。

本文档交付对象为项目经理、业务方、教师与学生用户，作为系统功能验收与使用培训的核心依据。

本文档的读者对象包括：

\begin{itemize}
	\item 项目管理人员：了解系统整体功能范围与交付边界，用于项目规划与进度管控；
	\item 业务方/甲方代表：确认系统功能是否满足实际业务需求，作为验收依据；
	\item 测试人员：依据功能需求条目编写测试用例，验证系统功能完整性与正确性；
	\item 最终用户（教师、学生、管理员）：了解系统提供的功能与操作方式，作为使用参考。
\end{itemize}

\subsection{项目背景与范围}
本系统为高校软件工程实训项目，服务对象为高校教师与在校学生。

在当前高校教学实践中，学生撰写课程论文与毕业论文时普遍面临以下问题：
\begin{enumerate}
	\item 论文格式不统一，不同学院、不同课程的论文模板差异较大，学生排版耗时且容易出错；
	\item Word 排版工具门槛较高，部分学生不熟悉样式、分节符、页眉页脚等高级功能，导致格式反复修改；
	\item 教师批阅纸质或电子版论文时，批注留痕不系统，修改意见分散，难以追溯历史版本；
	\item 论文模板由各学院自行管理，缺乏统一的模板配置与分发机制，更新不及时。
\end{enumerate}

针对上述问题，本系统旨在提供一套基于 Web 的在线论文编辑与标准化文档生成工具，实现从模板配置、在线写作、格式预览、文档导出到教师批阅的完整业务闭环。

系统核心业务范围涵盖以下五个方面：
\begin{enumerate}
	\item \textbf{模板管理}：管理员对论文模板的创建、编辑、版本控制与上下架管理；
	\item \textbf{在线写作}：学生基于模板进行结构化论文撰写，支持图文表格与参考文献插入；
	\item \textbf{文档导出}：按模板格式一键导出标准 docx 与 PDF 格式文件；
	\item \textbf{在线预览}：Web 端整篇论文完整排版预览与分页浏览；
	\item \textbf{课程批阅}：教师对学生提交的课程论文进行批注、评分与版本管理。
\end{enumerate}

系统独立部署运行，无需对接第三方外部系统。

\subsection{定义与缩略语}
\begin{itemize}
	\item \textbf{富文本编辑器}：网页端所见即所得的文本编辑组件，支持字体、层级、表格、图片嵌入等排版操作；
	\item \textbf{docx}：Microsoft Word 文档标准格式；
	\item \textbf{PDF}：便携式文档格式，用于跨平台保持文档排版一致性；
	\item \textbf{三线表}：学术论文常用表格格式，由顶线、栏目线和底线三条横线构成；
	\item \textbf{模板}：论文的标准格式框架，包含章节结构、排版参数与封面字段定义，由管理员创建并维护；
	\item \textbf{草稿}：学生正在编辑中、尚未提交的论文状态，系统自动保存以防止数据丢失；
	\item \textbf{批注}：教师在论文特定位置添加的修改意见或建议，学生可见并据此修改。
\end{itemize}

\subsection{参考资料}
\begin{enumerate}
	\item GB/T 9385-2008 计算机软件需求规格说明规范；
	\item GB/T 7714-2015 信息与文献 参考文献著录规则；
	\item 实训项目任务书：在线论文文档自动生成工具开发（2026 软件工程实训）。
\end{enumerate}

\subsection{文档结构说明}
本文档按照以下结构组织：
\begin{enumerate}
	\item \textbf{引言}：说明编写目的、项目背景、术语定义与参考资料；
	\item \textbf{任务概述}：描述系统目标、用户角色、使用环境与约束条件；
	\item \textbf{功能需求}：按模块详细描述系统必须具备的各项功能；
	\item \textbf{数据需求}：描述系统涉及的核心数据实体与业务规则；
	\item \textbf{非功能性需求}：描述性能、安全性、易用性、可靠性等方面的要求；
	\item \textbf{运行需求}：描述用户界面、故障处理等方面的要求；
	\item \textbf{附录}：提供模板结构示例与参考文献格式说明。
\end{enumerate}

\section{任务概述}

\subsection{系统目标}
开发一套 Web 端在线论文编辑与标准化文档生成系统，主要用于辅助学生规范化完成课程论文、毕业论文、项目设计类论文的在线撰写、版式排版与文档导出。

系统设置管理员、学生、教师三类用户角色，分别赋予模板管理、论文编辑导出、课程论文线上批阅等对应权限，实现论文模板统一配置、网页端图文表格参考文献编辑、全文在线预览、docx 与 PDF 格式文件导出、课程论文批阅留痕等完整业务流程。

系统预期实现以下业务目标：
\begin{enumerate}
	\item 学生论文排版格式统一率达到 100\%，消除因手动排版导致的格式不一致问题；
	\item 学生论文排版时间减少 50\% 以上，系统自动套用模板格式，学生专注于内容写作；
	\item 教师批阅效率提升，批注意见系统化留存，支持退回修改与复审闭环；
	\item 模板配置标准化，各学院论文模板统一由管理员维护，版本可追溯。
\end{enumerate}

\subsection{系统角色与权限划分}
系统面向三类用户，各自拥有独立的操作权限，互不越界：

\subsubsection{系统管理员}
管理员是系统后台的管理者，负责基础数据与模板的维护工作。具体权限包括：
\begin{itemize}
	\item 学院信息管理：新增、编辑、删除学院信息，按名称搜索学院；
	\item 模板管理：新增论文模板，配置模板的封面字段、章节结构与排版格式；
	\item 模板编辑：修改已有模板的各项配置，版本号自动递增；
	\item 模板上下架：启用或停用模板，控制模板对学生端的可见性；
	\item 模板删除：删除不再使用的模板；
	\item 模板详情查看：以只读方式查看模板的完整配置信息。
\end{itemize}

管理员不可直接编辑学生论文内容，不可参与论文批阅流程。

\subsubsection{学生用户}
学生是系统的主要使用群体，通过系统完成论文的在线撰写与导出。具体权限包括：
\begin{itemize}
	\item 模板选用：按学院与论文类型筛选可用模板，选择后创建论文项目；
	\item 论文编辑：在富文本编辑器中分章节撰写论文内容，支持文字排版、图片上传、表格插入；
	\item 参考文献管理：逐条录入参考文献信息，在正文中插入引用标注；
	\item 草稿保存：系统自动保存编辑内容，学生也可手动保存；
	\item 全文预览：一键生成整篇论文的网页预览，查看排版效果；
	\item 文档导出：按模板格式导出 docx 或 PDF 文件；
	\item 课程论文提交：完成编辑后将课程论文提交给授课教师；
	\item 查看批阅结果：查看教师批注、评语与评分，根据意见修改后重新提交。
\end{itemize}

学生仅可编辑自己创建的论文，不可查看或修改其他学生的论文。

\subsubsection{教师用户}
教师是课程论文批阅的执行者，仅在课程论文场景下拥有操作权限。具体权限包括：
\begin{itemize}
	\item 查看待批阅列表：查看学生提交给自己的课程论文列表；
	\item 论文预览：在线查看学生提交的论文内容与排版效果；
	\item 添加批注：在论文任意段落或位置添加批注与修改建议；
	\item 评阅打分：填写整体评语，评定论文分数与等级；
	\item 退回论文：将不符合要求的论文退回给学生，需填写退回原因与修改建议；
	\item 通过论文：对符合要求的论文执行通过操作；
	\item 查看批阅历史：查看某篇论文历次提交、批阅、退回的完整记录。
\end{itemize}

教师仅可批阅提交给自己的课程论文，不可批阅其他教师的论文，不可编辑论文内容。

\subsection{用户特征描述}
\begin{itemize}
	\item \textbf{管理员}：具备基本计算机操作能力，熟悉 Web 后台管理界面操作，通常为教务管理人员或实训指导教师；
	\item \textbf{学生}：具备基本计算机操作能力，熟悉 Word 等文字处理软件，但可能不熟悉高级排版功能；使用频率集中在课程论文与毕业设计期间；
	\item \textbf{教师}：具备较高计算机操作能力，主要使用系统的批阅功能，关注批注的便捷性与评阅效率。
\end{itemize}

\subsection{系统使用环境}
\begin{itemize}
	\item 客户端：主流浏览器（Chrome、Edge、Firefox 最新版本）；
	\item 网络环境：支持 HTTP/HTTPS 协议的校园网或互联网环境；
	\item 显示设备：PC 端显示器或笔记本电脑，推荐分辨率 1920$\times$1080 及以上。
\end{itemize}

系统需满足以下使用约束：
\begin{enumerate}
	\item 单篇论文字数限制不超过 50000 字；
	\item 单篇论文图片不超过 50 张，单张图片大小不超过 5MB；
	\item 参考文献不超过 200 条；
	\item 教师批注单条不超过 500 字符。
\end{enumerate}

\subsection{系统边界与限制}
本系统\textbf{不}具备以下功能，超出范围的需求不在本次开发目标内：
\begin{enumerate}
	\item 不提供论文查重功能，系统不检测论文内容的原创性；
	\item 不提供 AI 写作辅助功能，论文内容完全由学生自主撰写；
	\item 不提供学术不端检测功能，抄袭判定由教师或学校其他系统负责；
	\item 不对接学校教务系统，课程信息与学生名单不自动同步；
	\item 不提供移动端适配，系统仅支持 PC 端浏览器访问；
	\item 不提供论文内容的版本差异对比功能（如 Word 的修订模式）。
\end{enumerate}

\section{功能需求}

\subsection{功能划分}
系统划分为以下六大核心功能模块：
\begin{enumerate}
	\item 用户认证与权限管理模块；
	\item 后台基础数据管理模块（学院管理）；
	\item 论文模板管理模块；
	\item 学生在线写作模块；
	\item 文档预览与导出模块；
	\item 课程论文批阅模块。
\end{enumerate}

各模块之间的业务关系如下：
\begin{itemize}
	\item 模板管理模块为在线写作模块提供论文格式规范，学生创建论文时必须选用一个已启用的模板；
	\item 在线写作模块为预览与导出模块提供论文内容数据，预览与导出均基于当前论文内容生成；
	\item 在线写作模块与课程论文批阅模块存在提交-批阅流转关系，学生提交后教师方可批阅；
	\item 学院管理模块为模板管理模块提供学院分类数据，模板必须归属某个学院。
\end{itemize}

\subsection{功能需求详细描述}

\subsubsection{用户认证与权限管理}

该模块负责用户的注册、登录与权限隔离，是系统安全运行的基础保障。

\reqitem{FR001}{系统支持用户通过用户名和密码自行注册账号，需填写用户名、密码及角色信息（管理员/学生/教师），系统校验用户名唯一性，重复用户名不允许注册。密码采用加密方式存储，系统内部不保留明文密码。}

\reqitem{FR002}{系统支持所有已注册用户通过用户名和密码登录。登录成功后，系统根据用户角色自动跳转至对应的功能界面：管理员进入后台管理界面，学生进入论文工作台，教师进入批阅工作台。登录失败时，系统提示"用户名或密码错误"。}

\reqitem{FR003}{系统对三类角色进行严格访问隔离。管理员端功能仅管理员可见和操作，学生端功能仅学生可见和操作，教师端功能仅教师可见和操作。任何用户尝试访问非本角色的功能区域时，系统拒绝访问并提示"权限不足"。}

\reqitem{FR004}{用户可随时通过界面上的"退出登录"按钮退出系统。退出后，用户需重新输入账号密码才能再次进入系统。退出时系统清除当前登录状态，确保他人无法继续使用该账号。}

\textbf{业务规则：}
\begin{itemize}
	\item 用户名长度不超过 50 个字符，不可包含空格和特殊字符；
	\item 密码长度不少于 6 个字符；
	\item 已被禁用的账号无法登录，系统提示"账户已被禁用"。
\end{itemize}

\subsubsection{后台基础数据管理（学院管理）}

该模块由管理员操作，维护学院基础信息，为模板分类提供数据支撑。

\reqitem{FR005}{管理员可查看所有学院列表，列表以表格形式展示学院名称、描述、创建时间等信息，支持按学院名称关键字进行模糊搜索，搜索结果实时更新。}

\reqitem{FR006}{管理员可新增学院信息，填写学院名称（必填，不可与已有学院重复）和描述信息（选填，简要说明该学院）。新增成功后，学院立即出现在列表中。}

\reqitem{FR007}{管理员可编辑已有学院的名称和描述信息。编辑后保存即可生效，已关联该学院的模板不受影响。}

\reqitem{FR008}{管理员可删除学院，删除前系统弹出二次确认对话框，明确提示"删除后无法恢复"。确认删除后，学院从列表中移除。}

\textbf{业务规则：}
\begin{itemize}
	\item 学院名称不能为空，长度不超过 100 个字符；
	\item 同一系统内学院名称不可重复；
	\item 学院描述长度不超过 500 个字符；
	\item 删除学院前，管理员应确认该学院下没有关联的论文模板，避免数据引用异常。
\end{itemize}

\subsubsection{论文模板管理}

该模块是管理员端的核心功能，负责论文模板的全生命周期管理。模板定义了论文的格式规范，学生创建论文时自动套用所选模板的格式。

\reqitem{FR009}{管理员可查看所有模板列表，列表以表格形式展示模板名称、类型、所属学院、启用状态、版本号、更新时间等信息。支持按学院、模板类型（毕业论文/课程论文/项目论文）、启用状态（已启用/已停用）进行组合筛选，搜索结果实时更新。}

\reqitem{FR010}{管理员可新增论文模板，填写基本信息并配置模板内容。新增时需完成以下配置：}

\textbf{基本信息：}
\begin{itemize}
	\item 模板名称：必填，长度不超过 200 个字符，用于在模板列表中标识模板；
	\item 模板类型：必填，从"毕业论文""课程论文""项目论文"三个选项中选择；
	\item 所属学院：必填，从已有的学院列表中选择，模板将归属于该学院分类下；
	\item 描述信息：选填，简要说明模板用途或适用场景；
	\item 启用状态：默认为"停用"状态，管理员可在新增时选择直接启用。
\end{itemize}

\textbf{封面字段配置：}
管理员可定义论文封面上需要学生填写的字段。每个字段包含以下属性：
\begin{itemize}
	\item 字段标签：显示在封面上的名称（如"论文题目""学院""专业"等）；
	\item 字段标识：字段的英文标识符（如"title""college""major"等），用于系统内部引用；
	\item 字段类型：文本类型或日期类型；
	\item 是否必填：控制该字段是否为学生必须填写的项目；
	\item 最大长度：限制学生输入内容的最大字符数。
\end{itemize}
系统预设默认封面字段包括：论文题目、学院、专业、学生姓名、学号、指导老师、指导老师职称、完成日期。管理员可根据需要新增、删除字段，或调整字段排列顺序。

\textbf{章节结构配置：}
管理员可配置论文包含的章节及其属性。系统预设以下默认章节：
\begin{enumerate}
	\item 封面（必填、显示、不可编辑——由封面字段配置自动渲染）
	\item 原创声明（必填、显示、不可编辑——固定模板文本）
	\item 中文摘要（必填、显示、可编辑）
	\item 中文关键词（必填、显示、可编辑——关键词之间用分号分隔）
	\item 英文摘要（必填、显示、可编辑）
	\item 英文关键词（必填、显示、可编辑）
	\item 目录（必填、显示、不可编辑——系统根据章节标题自动生成）
	\item 正文（必填、显示、可编辑——支持最多三级标题，学生可自由增删章节）
	\item 参考文献（必填、显示、可编辑——由参考文献模块单独管理）
	\item 致谢（非必填、显示、可编辑——学生可选填）
	\item 附录（非必填、显示、可编辑——学生可选填）
\end{enumerate}
管理员可对每个章节配置以下属性：
\begin{itemize}
	\item 是否必填：控制该章节是否为学生必须完成的部分；
	\item 是否显示：控制该章节是否出现在论文中（设为不显示则整节隐藏）；
	\item 是否允许编辑：控制学生是否可以直接编辑该节内容（封面、原创声明、目录等固定内容不可编辑）。
\end{itemize}

\textbf{排版格式配置：}
管理员可配置论文各组成部分的排版参数，具体包括：
\begin{itemize}
	\item \textbf{页面设置}：页面尺寸（A4/A3/B5）、上下左右页边距（单位：厘米）；
	\item \textbf{正文样式}：字体、字号、行间距、首行缩进、对齐方式；
	\item \textbf{一级标题}：字体、字号、粗体、对齐方式、段前段后间距、是否自动编号；
	\item \textbf{二级标题}：同上；
	\item \textbf{三级标题}：同上；
	\item \textbf{摘要样式}：标题字体字号与对齐、正文字体字号与行距；
	\item \textbf{参考文献}：标题样式（字体、字号、对齐）、正文样式（字体、字号、行距）、编号格式（默认为 [N] 格式）、是否悬挂缩进；
	\item \textbf{致谢样式}：标题样式（字体、字号、对齐）、正文样式（字体、字号、行距、首行缩进）；
	\item \textbf{封面标题样式}：字体、字号、粗体、对齐方式；
	\item \textbf{封面正文样式}：字体、字号、对齐方式、行距；
	\item \textbf{页眉}：页眉文字（支持学院名称占位符，渲染时自动替换）、字体、字号、对齐方式、是否显示分隔线；
	\item \textbf{页脚}：是否显示页码、页码格式（阿拉伯数字/罗马数字）、字体、字号、对齐方式。
\end{itemize}

新增成功后，模板默认为版本 1，状态为停用（除非管理员在新增时选择了启用）。

\reqitem{FR011}{管理员可编辑已有模板的基本信息、封面字段、章节结构和排版格式。每次编辑保存后，模板版本号自动递增（如从 v1 升级为 v2）。模板更新后，学生新建论文将自动套用最新版本的模板格式，已创建的论文不受影响。}

\reqitem{FR012}{管理员可启用或停用模板。停用后，该模板在学生端的模板列表中不再显示，学生新建论文时不可选用该模板。已使用该模板创建的论文不受影响，学生可继续编辑和导出。}

\reqitem{FR013}{管理员可删除模板，删除前系统弹出二次确认对话框。删除后模板不可恢复，但已使用该模板创建的论文不受影响，论文中保留创建时的模板格式快照。}

\reqitem{FR014}{管理员可查看模板详情，以只读方式展示模板的全部配置信息，包括基本信息、章节结构表格、封面字段表格和排版格式参数。详情页不提供编辑功能。}

\reqitem{FR015}{模板更新后，学生新建论文将自动套用最新版本的模板格式。已创建的论文不受模板更新影响，继续使用创建时的模板格式。}

\textbf{业务规则：}
\begin{itemize}
	\item 模板名称不能为空，长度不超过 200 个字符；
	\item 模板必须归属某个学院，不可创建无归属的模板；
	\item 模板类型必须为"毕业论文""课程论文""项目论文"三者之一；
	\item 每次编辑保存，版本号自动 +1，不可手动指定版本号；
	\item 删除模板时，已使用该模板的论文数据不受影响。
\end{itemize}

\subsubsection{学生在线写作}

该模块是学生端的核心功能，学生在此完成论文的全部编辑工作。

\reqitem{FR016}{学生可按学院与论文类型筛选可用模板，从已启用的模板中选择一个，点击创建论文项目。创建时系统自动根据所选模板生成预设的章节结构，学生即可开始编辑。}

\reqitem{FR017}{编辑器支持以下文字排版操作：标题层级设置（一级/二级/三级标题）、文字加粗、文字斜体、下划线、左对齐/居中/右对齐/两端对齐、段落缩进、有序列表/无序列表。编辑器为所见即所得模式，学生可直接在页面上编辑内容。}

\reqitem{FR018}{学生可从本地上传图片（支持 JPG、PNG、GIF、SVG 格式），上传后图片嵌入文档指定位置。单张图片大小不超过 5MB，单篇论文图片总数不超过 50 张。图片上传失败时，系统提示具体原因（如格式不支持、文件过大等）。}

\reqitem{FR019}{学生可在文档中插入自定义表格，支持设置行列数、合并拆分单元格、调整列宽行高。系统提供学术三线表模板，学生可一键将普通表格转换为三线表格式，满足学术论文的表格规范要求。}

\reqitem{FR020}{学生可通过独立的参考文献管理模块逐条录入文献信息，每条文献包含以下字段：文献类型（期刊论文/专著/学位论文/会议论文/网页等）、作者、标题、刊物名称、出版年份、卷号、期号、页码。系统按照模板定义的参考文献格式（默认 [N] 编号格式）自动统一排版参考文献列表。}

\reqitem{FR021}{学生可在正文中插入引用标注，系统自动生成带方括号的序号（如 [1]、[2]），并与参考文献条目关联。引用标注的编号随参考文献列表的顺序自动更新。}

\reqitem{FR022}{系统每 30 秒自动保存一次论文草稿，同时学生可随时点击"保存"按钮手动保存。自动保存在后台静默执行，不影响学生的编辑操作。学生可多次退出并重新进入系统，上次保存的内容会自动加载，支持续写修改。}

\textbf{业务规则：}
\begin{itemize}
	\item 学生仅可编辑自己创建的论文，不可查看或修改其他学生的论文；
	\item 已提交给教师的课程论文处于锁定状态，学生不可继续编辑，除非教师将其退回；
	\item 图片上传仅支持 JPG、PNG、GIF、SVG 格式，单张不超过 5MB；
	\item 参考文献条目不超过 200 条；
	\item 论文总字数不超过 50000 字。
\end{itemize}

\subsubsection{全文预览与文档导出}

该模块为学生提供论文的最终效果查看与文件下载功能。

\reqitem{FR023}{学生可一键生成整篇论文的网页预览页面。预览页面完整展示论文的封面、各级章节、图片、表格、参考文献的排版效果，模拟 A4 纸张样式，支持分页浏览。预览页面为只读模式，不可编辑。}

\reqitem{FR024}{系统按模板既定格式，一键导出标准 docx 格式 Word 文档。导出的 Word 文档完整保留模板定义的排版格式（字体、字号、行距、页边距、页眉页脚等），图片与表格正确嵌入，参考文献按模板格式统一排版。同时支持导出 PDF 格式文件，PDF 排版与 docx 保持一致。}

\reqitem{FR025}{导出功能支持选择导出范围：可选择导出整篇论文（包含封面、摘要、正文、参考文献等全部章节），或仅导出正文部分（从第一章开始到最后一章结束）。}

\textbf{业务规则：}
\begin{itemize}
	\item 预览与导出基于最近一次保存的论文内容，未保存的编辑不会体现在预览和导出结果中；
	\item 导出的文件名默认为论文标题，文件名中不可包含特殊字符；
	\item 空章节在导出时自动跳过，不生成空白页面。
\end{itemize}

\subsubsection{课程论文批阅}

该模块实现学生提交课程论文后，教师在线批阅、评分、退回的完整业务流程。

\reqitem{FR026}{学生完成课程论文编辑后，可将论文提交给对应授课教师。提交后，该论文进入锁定状态，学生不可继续编辑。提交操作不可撤销，学生需确认后再执行。}

\reqitem{FR027}{教师可在线查看学生提交的论文预览内容，完整浏览论文排版效果。教师可在论文的任意段落或行位置添加批注，批注内容为修改建议或问题说明。每条批注不超过 500 个字符。}

\reqitem{FR028}{教师可填写整体评语，评定论文分数（0-100 分）与等级（优秀/良好/中等/及格/不及格）。教师执行"通过"操作表示论文合格，执行"退回"操作时必须填写退回原因与修改建议，退回后学生可继续修改。}

\reqitem{FR029}{系统留存每一次提交、批阅、退回的版本记录。学生可查看教师添加的所有批注内容、评语与评分，根据修改意见调整论文后可重新提交复审。每次重新提交，教师可查看此前所有批阅历史，了解论文的修改过程。}

\textbf{业务规则：}
\begin{itemize}
	\item 仅课程论文类型支持提交与批阅流程，毕业论文和项目论文不参与此流程；
	\item 论文提交后锁定，学生不可编辑，除非教师执行退回操作；
	\item 退回后学生可修改并重新提交，提交次数不限；
	\item 教师仅可批阅提交给自己的论文，不可批阅其他教师的论文；
	\item 评分范围为 0-100 分，等级与分数对应关系为：优秀（90-100）、良好（80-89）、中等（70-79）、及格（60-69）、不及格（0-59）；
	\item 批阅历史永久保留，不可删除。
\end{itemize}

\subsection{功能需求汇总表}

\begin{tabular}{|p{1.5cm}|p{6cm}|p{2.5cm}|}
	\hline
	\textbf{编号} & \textbf{功能描述} & \textbf{所属模块} \\
	\hline
	FR001 & 用户注册（用户名+密码+角色） & 认证管理 \\
	\hline
	FR002 & 用户登录与角色跳转 & 认证管理 \\
	\hline
	FR003 & 角色访问隔离与权限控制 & 认证管理 \\
	\hline
	FR004 & 退出登录与状态清除 & 认证管理 \\
	\hline
	FR005 & 学院列表查看与搜索 & 学院管理 \\
	\hline
	FR006 & 新增学院 & 学院管理 \\
	\hline
	FR007 & 编辑学院信息 & 学院管理 \\
	\hline
	FR008 & 删除学院 & 学院管理 \\
	\hline
	FR009 & 模板列表查看与筛选 & 模板管理 \\
	\hline
	FR010 & 新增模板（含封面/章节/格式配置） & 模板管理 \\
	\hline
	FR011 & 编辑模板配置 & 模板管理 \\
	\hline
	FR012 & 模板启用/停用 & 模板管理 \\
	\hline
	FR013 & 删除模板 & 模板管理 \\
	\hline
	FR014 & 查看模板详情（只读） & 模板管理 \\
	\hline
	FR015 & 模板更新自动套用 & 模板管理 \\
	\hline
	FR016 & 选用模板创建论文 & 在线写作 \\
	\hline
	FR017 & 富文本编辑（排版操作） & 在线写作 \\
	\hline
	FR018 & 图片上传与嵌入 & 在线写作 \\
	\hline
	FR019 & 表格插入与三线表转换 & 在线写作 \\
	\hline
	FR020 & 参考文献录入与排版 & 在线写作 \\
	\hline
	FR021 & 正文引用标注 & 在线写作 \\
	\hline
	FR022 & 自动保存与手动保存 & 在线写作 \\
	\hline
	FR023 & 全文网页预览 & 预览导出 \\
	\hline
	FR024 & docx/PDF 文档导出 & 预览导出 \\
	\hline
	FR025 & 导出范围选择 & 预览导出 \\
	\hline
	FR026 & 课程论文提交与锁定 & 论文批阅 \\
	\hline
	FR027 & 教师添加批注 & 论文批阅 \\
	\hline
	FR028 & 评阅打分与退回 & 论文批阅 \\
	\hline
	FR029 & 批阅历史查看与重新提交 & 论文批阅 \\
	\hline
\end{tabular}

\section{数据需求}

\subsection{核心数据实体}
系统涉及以下核心数据实体：

\subsubsection{用户}
用户是系统的操作主体，包含以下关键信息：用户名（唯一标识，用于登录）、密码（加密存储）、昵称、角色（管理员/学生/教师）、所属学院、邮箱、账号状态（正常/禁用）、注册时间、最后更新时间。

\subsubsection{学院}
学院是模板分类的基础维度，包含以下关键信息：学院名称（唯一）、描述信息、创建时间、更新时间。

\subsubsection{论文模板}
论文模板定义了论文的格式规范，包含以下关键信息：模板名称、模板类型（毕业论文/课程论文/项目论文）、所属学院 ID、启用状态（启用/停用）、版本号（每次编辑自动递增）、描述信息、创建时间、更新时间。

\subsubsection{模板配置}
模板配置存储模板的详细格式定义，与模板一一关联，包含以下关键信息：
\begin{itemize}
	\item 章节结构（JSON 格式）：存储各章节的名称、是否必填、是否显示、是否允许编辑、排列顺序等信息；
	\item 排版格式（JSON 格式）：存储页面、正文、各级标题、摘要、参考文献、致谢、封面、页眉页脚的排版参数；
	\item 封面字段（JSON 格式）：存储封面各字段的标签名、标识符、类型、是否必填、最大长度等信息。
\end{itemize}

\subsubsection{论文}
论文是学生创建的文档实体，包含以下关键信息：论文标题、学生 ID、论文内容（长文本，存储富文本 HTML）、论文状态（草稿/已提交/已批阅/已退回）、创建时间、更新时间。

\subsection{数据状态流转}

\subsubsection{论文状态流转}
论文在生命周期中经历以下状态变化：
\begin{enumerate}
	\item \textbf{草稿}：学生创建论文后的初始状态，学生可自由编辑、保存；
	\item \textbf{已提交}：学生将课程论文提交给教师后进入此状态，论文锁定，学生不可编辑；
	\item \textbf{已批阅}：教师执行"通过"操作后进入此状态，论文锁定，学生可查看评阅结果；
	\item \textbf{已退回}：教师执行"退回"操作后进入此状态，论文解锁，学生可继续编辑并重新提交。
\end{enumerate}

状态流转规则：
\begin{itemize}
	\item 草稿 $\rightarrow$ 已提交：学生主动提交；
	\item 已提交 $\rightarrow$ 已批阅：教师执行通过操作；
	\item 已提交 $\rightarrow$ 已退回：教师执行退回操作；
	\item 已退回 $\rightarrow$ 已提交：学生修改后重新提交。
\end{itemize}

\subsubsection{模板状态流转}
模板在生命周期中经历以下状态变化：
\begin{enumerate}
	\item \textbf{已启用}：模板对学生端可见，学生可选用该模板创建论文；
	\item \textbf{已停用}：模板对学生端不可见，学生新建论文时无法选用。
\end{enumerate}

状态流转规则：
\begin{itemize}
	\item 已启用 $\leftrightarrow$ 已停用：管理员随时切换，已使用该模板的论文不受影响。
\end{itemize}

\subsection{业务规则汇总}

\subsubsection{注册与登录规则}
\begin{enumerate}
	\item 用户名不可重复，长度不超过 50 个字符；
	\item 密码长度不少于 6 个字符，加密存储；
	\item 已禁用账号无法登录。
\end{enumerate}

\subsubsection{学院管理规则}
\begin{enumerate}
	\item 学院名称不可重复，长度不超过 100 个字符；
	\item 学院描述长度不超过 500 个字符；
	\item 删除学院前应确认无关联模板。
\end{enumerate}

\subsubsection{模板管理规则}
\begin{enumerate}
	\item 模板必须归属某个学院；
	\item 模板类型必须为"毕业论文""课程论文""项目论文"之一；
	\item 每次编辑保存版本号自动 +1；
	\item 停用的模板不影响已创建的论文；
	\item 删除的模板不影响已创建的论文。
\end{enumerate}

\subsubsection{论文编辑规则}
\begin{enumerate}
	\item 学生仅可编辑自己创建的论文；
	\item 已提交的论文处于锁定状态，不可编辑；
	\item 退回后的论文可继续编辑；
	\item 图片单张不超过 5MB，总数不超过 50 张；
	\item 参考文献不超过 200 条；
	\item 论文总字数不超过 50000 字。
\end{enumerate}

\subsubsection{批阅规则}
\begin{enumerate}
	\item 仅课程论文支持提交与批阅流程；
	\item 教师仅可批阅提交给自己的论文；
	\item 评分范围 0-100 分；
	\item 退回时必须填写修改建议；
	\item 批阅历史永久保留。
\end{enumerate}

\section{非功能性需求}

\subsection{性能需求}
\begin{itemize}
	\item 页面加载响应时间不超过 2 秒（标准网络环境）；
	\item 论文保存操作响应时间不超过 1 秒；
	\item 网页预览生成时间不超过 3 秒；
	\item 文档导出时间不超过 10 秒（单篇论文万字以内）；
	\item 参考文献序号与正文引用标注一一对应，准确无误；
	\item 导出文档排版格式与模板定义一致。
\end{itemize}

\subsection{安全性需求}
\begin{itemize}
	\item \textbf{密码安全}：用户密码采用加密存储，严禁明文保存；
	\item \textbf{越权防范}：学生仅可编辑本人论文，教师仅可批阅本人课程论文，管理员不可编辑学生论文内容；
	\item \textbf{数据保护}：用户个人信息与论文内容受系统保护，非授权用户不可访问。
\end{itemize}

\subsection{易用性需求}
\begin{itemize}
	\item 界面简洁清晰，操作反馈明确；
	\item 重要操作（删除、提交、退出）需二次确认弹窗；
	\item 表单填写提供必填项提示与输入校验，填写不合法时给出明确错误提示；
	\item 操作结果通过 Toast 提示反馈（成功显示绿色、失败显示红色、警告显示黄色）；
	\item 支持主流屏幕分辨率（1920$\times$1080、1366$\times$768）。
\end{itemize}

\subsection{可靠性需求}
\begin{itemize}
	\item 自动保存机制确保最多丢失 30 秒内编辑内容；
	\item 系统具备故障恢复能力，已保存数据不丢失；
	\item 网络中断时，前端提示用户当前网络状态，避免用户在断网环境下继续编辑导致数据丢失。
\end{itemize}

\subsection{可维护性需求}
\begin{itemize}
	\item 系统采用前后端分离架构，前端与后端独立开发与部署；
	\item 模板配置采用 JSON 格式存储，便于扩展新的章节类型与排版参数；
	\item 系统界面采用组件化设计，新增页面或功能模块时可复用现有组件。
\end{itemize}

\section{运行需求}

\subsection{用户界面}
系统采用 Web 可视化界面，整体 UI 风格以浅蓝色、白色、浅灰色为主色调，简洁清晰：

\paragraph{登录/注册页面}
居中卡片式布局，包含用户名输入框、密码输入框、角色选择（注册时）、登录/注册按钮。登录失败时在页面上方显示错误提示。

\paragraph{管理员后台界面}
采用左侧导航菜单 + 右侧内容区的布局：
\begin{itemize}
	\item 左侧导航：包含"学院管理"和"模板列表"两个菜单项，当前选中项高亮显示；
	\item 顶部栏：显示面包屑导航（管理后台 / 当前页面名称）和当前登录用户名；
	\item 内容区：以卡片形式承载页面内容，表格化展示数据，提供搜索框、筛选条件和操作按钮；
	\item 弹窗：新增/编辑操作通过模态弹窗完成，弹窗包含表单输入区域和"取消""确定"按钮；
	\item 确认弹窗：删除等危险操作前弹出二次确认对话框，明确提示操作后果。
\end{itemize}

\paragraph{模板编辑界面}
模板编辑采用多标签页布局，分四个标签页组织：
\begin{enumerate}
	\item \textbf{基本信息}：模板名称、类型、所属学院、状态、描述的表单输入；
	\item \textbf{封面字段}：可视化列表编辑器，支持新增、删除、上下移动字段，每个字段行内编辑标签名、标识符、类型、长度、是否必填；
	\item \textbf{章节结构}：表格形式展示各章节，每行提供"必填""显示""可编辑"三个复选框，正文章节额外提供最大层级设置；
	\item \textbf{格式参数}：分组表单，按页面设置、正文样式、各级标题、摘要、参考文献、致谢、封面、页眉页脚分组展示排版参数。
\end{enumerate}

\paragraph{学生工作台}
左侧导航菜单，右侧内容展示区，包含论文列表（展示论文标题、状态、创建时间、更新时间）和操作按钮。

\paragraph{论文编辑器}
顶部工具栏（排版操作按钮），中部富文本编辑区（按章节分区域编辑），底部状态栏（显示保存状态）与保存按钮。

\paragraph{论文预览页}
模拟 A4 纸张样式，完整展示论文排版效果，支持分页浏览，提供导出按钮。

\paragraph{教师批阅界面}
左侧论文预览区（展示学生提交的论文），右侧批注面板（添加批注），底部评语与评分表单。

\subsection{故障处理}
\begin{itemize}
	\item \textbf{网络请求失败}：弹窗提示错误原因，保留用户当前编辑内容，用户可关闭弹窗后重试；
	\item \textbf{文件上传失败}：弹窗提示具体原因（格式不支持、文件过大、网络超时等），保留用户编辑内容；
	\item \textbf{文档导出超时}：提示"导出任务超时，请稍后重试"；
	\item \textbf{网络中断}：前端提示"网络已断开"，恢复后用户可手动刷新页面重试；
	\item \textbf{服务器异常}：提示"服务器暂时不可用，请稍后重试"。
\end{itemize}

\section{附录}

\subsection{附录A：论文模板默认章节结构}

以本科毕业论文模板为例，预设章节结构如下：

\begin{enumerate}
	\item 封面（自动生成：标题、姓名、学号、学院、专业、指导教师、提交日期）
	\item 原创声明（固定模板，学生签名确认）
	\item 中文摘要与关键词（学生填写，关键词以分号分隔）
	\item 英文摘要与关键词（对应中文翻译）
	\item 目录（系统根据章节标题自动生成，不可手动编辑）
	\item 正文（支持最多三级标题，学生可自由增删章节）
	\item 参考文献（由参考文献模块单独管理，编号格式为 [N]）
	\item 致谢（可选）
	\item 附录（可选）
\end{enumerate}

\subsection{附录B：参考文献格式示例}

参考文献格式遵循 GB/T 7714-2015 国家标准，默认编号格式为 [N]，示例如下：

\textbf{期刊论文：}
[1] 张三, 李四. 论文标题[J]. 期刊名称, 2024, 12(3): 45-52.

\textbf{专著：}
[2] 王五. 书名[M]. 北京: 出版社名称, 2023.

\textbf{学位论文：}
[3] 赵六. 论文题目[D]. 北京: 学校名称, 2024.

\textbf{会议论文：}
[4] 作者. 论文标题[C]// 会议名称. 出版地: 出版者, 年份: 页码.

\textbf{网页：}
[5] 作者. 文章标题[EB/OL]. (发布日期)[引用日期]. 网址.

\subsection{附录C：模板封面默认字段定义}

\begin{tabular}{|p{2.5cm}|p{2cm}|p{1.5cm}|p{1.2cm}|p{1.5cm}|}
	\hline
	\textbf{字段标签} & \textbf{字段标识} & \textbf{类型} & \textbf{必填} & \textbf{最大长度} \\
	\hline
	论文题目 & title & 文本 & 是 & 100 \\
	\hline
	学院 & college & 文本 & 是 & 50 \\
	\hline
	专业 & major & 文本 & 是 & 50 \\
	\hline
	学生姓名 & studentName & 文本 & 是 & 20 \\
	\hline
	学号 & studentId & 文本 & 是 & 20 \\
	\hline
	指导老师 & supervisor & 文本 & 是 & 20 \\
	\hline
	指导老师职称 & supervisorTitle & 文本 & 否 & 20 \\
	\hline
	完成日期 & date & 日期 & 是 & -- \\
	\hline
\end{tabular}

\subsection{附录D：排版格式默认参数}

以本科毕业论文模板的默认排版参数为例：

\textbf{页面设置：} A4 纸张，上下边距 2.54cm，左右边距 3.17cm。

\textbf{正文样式：} 宋体，12pt，1.5 倍行距，首行缩进 2 字符，两端对齐。

\textbf{一级标题：} 黑体，16pt，加粗，居中对齐，段前 24pt，段后 18pt，自动编号。

\textbf{二级标题：} 黑体，14pt，加粗，左对齐，段前 18pt，段后 12pt，自动编号。

\textbf{三级标题：} 黑体，12pt，加粗，左对齐，段前 12pt，段后 6pt，自动编号。

\textbf{摘要标题：} 黑体，16pt，居中对齐。

\textbf{摘要正文：} 宋体，12pt，1.5 倍行距。

\textbf{参考文献标题：} 黑体，16pt，居中对齐。

\textbf{参考文献正文：} 宋体，10.5pt，1.5 倍行距，编号格式 [N]，悬挂缩进 2 字符。

\textbf{致谢标题：} 黑体，16pt，居中对齐。

\textbf{致谢正文：} 宋体，12pt，1.5 倍行距，首行缩进 2 字符。

\textbf{封面标题：} 黑体，18pt，加粗，居中对齐。

\textbf{封面正文：} 宋体，14pt，左对齐，2 倍行距。

\textbf{页眉：} "XX大学本科毕业论文"，宋体，9pt，居中对齐，显示分隔线。

\textbf{页脚：} 显示页码（阿拉伯数字），Times New Roman，9pt，居中对齐。

\end{document}